\documentclass{../notatki}
\usepackage{physics}

\title{Algebra liniowa z zastosowaniami 2.\\ Podstawy obliczeń kwantowych}

\begin{document}

\section{Przestrzeń wektorowa}

Przestrzeń wektorowa $V$ nad ciałem $K$ to:
\begin{itemize}
  \item $(V, +)$ - grupa abelowa, $+$ - dodawanie wektorów
  \item $(K, \oplus, \odot)$ - ciało, $\oplus$ - dodawanie
    skalarów, $\odot$ - mnożenie skalarów
  \item $\cdot : K \times V \to V$ - mnożenie wektorów
  \item $
    \alpha \cdot (v_1 + v_2) = \alpha \cdot v_1 + \alpha \cdot v_2
    $
  \item $
    (\alpha_1 \oplus \alpha_2) \cdot v = \alpha_1 \cdot v + \alpha_2 \cdot v
    $
  \item $
    \alpha_1 \cdot (\alpha_2 \cdot v) = (\alpha_1 \odot \alpha_2) \cdot v
    $
  \item $
    1 \cdot v = v
    $
\end{itemize}

Intuicyjnie, jest to ciało skalarów i grupa wektorów z dodatkowym
działaniem. Przestrzenią wektorową może być wiele rzeczy. Od wielowymiarowych
ciał $\mathbb{R}^n$, poprzez bajty ($\mathbb{F}_2^8$), aż po funkcje w ramach
szeregów Fourier'a.

\subsection{Powłoka liniowa}

Dla $S = \{v_1, v_2, \dots, v_n\}$, czyli zbioru wektorów z $V$ nad $K$,
zbiór wszystkich kombinacji liniowych:
$$
\mathrm{span}(S) = \{ \alpha_1 v_1 + \alpha_2 v_2 + \dots + \alpha_n
v_n | \alpha_1, \alpha_2, \dots, \alpha_n \in K \}
$$
nazywamy powłoką zbioru $S$. Intuicyjnie, jest to zbiór skalarów odpowiadających
pewnym wektorom z $V$. Mnożymy skalary z wektorami i dostajemy nowy wektor,
powstały z kombinacji liniowej. Wszystkie takie wektory tworzą powłokę liniową.

Inaczej, powłoka liniowa to najmniejsza podprzestrzeń $V$, zawierająca wszystkie
wektory z $S$. Mówimy, że wektory z $S$ rozpinają podprzestrzeń $V$, równoważną
$\mathrm{span}(S)$.

\subsection{Liniowa niezależność}

Podzbiór $S$ z $V$ nad $K$ jest liniowo niezależny, jeśli dla dowolnego
podzbioru $S' \subseteq S = \{v_1, v_2, \dots, v_n\}$ i dowolnych
$\alpha_1, \alpha_2 \in K$:
$$
\alpha_1 v_1 + \alpha_2 v_2 + \dots + \alpha_n v_n  = 0 \rightarrow
\alpha_1 = \alpha_2 = \dots = \alpha_n = 0
$$

Intuicyjnie, jeśli nie da się wyrazić dowolnego wektora z $S$ jako kombinacji
liniowej innych wektorów z $S$, to $S$ jest liniowo niezależny. Jeśli da się
wyrazić dowolny wektor z $S$ przy pomocy innych wektorów z $S$, to da się też
zsumować dowolnie te wektory do zera, tak, że skalary nie są równe zero.

\subsection{Baza}

Baza $B$ przestrzeni $V$, to liniowo niezależny $B \subseteq V$, który rozpina
przestrzeń $V$.

Inaczej, jest to zbiór wektorów liniowo niezależnych, czyli takich których
nie da się wyrazić przy pomocy siebie, przy pomocy których można wyrazić dowolny
wektor z $V$ jako kombinację liniową. Można o bazie myśleć jako o zbiorze
elementarnych wartości, bloków budulcowych.

\subsection{Przekształcenia liniowe}

Niech $U$ i $V$ będą przestrzeniami wektorowymi nad $K$. Funkcję $T: U \to V$
nazywamy przekształceniem liniowym, jeśli dla dowolnych $\alpha_1,
\alpha_2 \in K$
i dowolnych $u_1, u_2 \in U$:
$$
T(\alpha_1 u_1 + \alpha_2 u_2) = \alpha_1 T(u_1) + \alpha_2 T(u_2)
$$

\subsection{Zmiana bazy}

Dla przestrzeni $U$ z bazą $S = \{v_1, v_2, \dots, v_n\}$, oraz $V$ z
bazą $S' = \{w_1, w_2, \dots, w_m\}$, macierzą przekształcenia liniowego
$T: U \to V$ nazywamy $M_S^{S'} = [m_{ij}] \in K^{m \times n}$.
$$
T(v_j) = \sum_{i=1}^m m_{ij} w_i
$$
Intuicyjnie, w j-tej kolumnie macierzy $M_S^{S'}$ stoją współrzędne wektora
$v_j$ w bazie $S'$.

Wspólrzędne wektora $v \in V$ w bazie $S'$ to $\langle v \rangle_{S'}
= M_S^{S'} \langle v \rangle_S$, dla danych koordynatów w bazie $S$
$\langle v \rangle_S$.

Macierzą przejścia dla przestrzeni $V$ z bazą S i $U$ z bazą $S'$,
oraz macierzą przekształcenia $M_S^{S'}$ jest macierz
$M_S^{S'}(Id_S)$, $Id_V : V \to V$.

\subsection{Własności}

Śladem macierzy $A$ nazywamy $\tr A$; gdzie:
$$
\tr A = \sum_{i=1}^n a_{ii}
$$

Jądrem nazywamy zbiór wektorów (punktów), które są przekształcane w zero.
$$
\ker T = \{v \in V: T(v) = 0\}
$$

Dla $T: V \to W$, obrazem jest podzbiór przestrzeni $W$, składający się z
wektorów, które można otrzymać z $T(v) : v \in V$.
$$
\operatorname{Im} T = T(V) = \{w \in W: \exists_{v \in V} T(v) = w\}
$$
Intuicyjnie, obrazem jest kontrdziedzina przekształcenia; wszystkie wartości
co może nam zwrócić.

$$
\dim V = \dim \operatorname{Ker} T + \dim \operatorname{Im} T
$$

\section{Wektory własne}

Dla przestrzeni $V$ nad ciałem $K$, z przekształceniem $T: V \to V$, jeśli
dla niezerowego $v \in V$ spełnione jest:
$$
Tv = \lambda v
$$
gdzie $\lambda \in K$, to $v$ nazywamy wektorem własnym dla opdowiadającej
wartości własnej $\lambda$. Tu się przydaje notacja Diraca;
$v_\lambda = \ket{\lambda}$. Intuicyjnie, wektor własny operatora (macierzy),
po wrzuceniu do operatora może zwrócić tylko przeskalowany ten sam wektor.
Odwrócenie macierzy nieosobliwej, odwraca też jej wartości własne.

Z powyższego można wyprowadzić wzór na wartości własne; oraz
zdefiniować wielomian charakterystyczny:
$$
\determinant (A - \lambda I) = 0 = p_T(\lambda)
$$
$$
(T - \lambda I)v_\lambda = 0
$$

Zbiór wszystkich wartości własnych $A$ nazywamy jej widmem (spektrum) i
oznaczamy $\sigma(A)$.

Wektor $\ket{m}$ nazywamy uogólnionym wektorem własnym rzędu $m$ macierzy $A$
dla wartości własnej $\lambda$ jeśli $(A - \lambda I)^m \ket{m} = 0$, oraz
$(A - \lambda I)^{m - 1} \ket{m} \neq 0$.
$$
\ket{n - 1} = (A - \lambda I) \ket{n}
$$

\subsection{Wielomian charakterystyczny}

Dla wielomianu charakterystycznego:
$$
p_A(t) = (t - \lambda_1)^{k_1} (t - \lambda_2)^{k_2} \dots (t - \lambda_m)^{k_m}
$$
liczbę $k_i : i \leq m$ nazywamy krotnością wartości własnej $\lambda_i$.
Jeśli da się zapisać $p_A(t) = c_nt^n + \dots c_0$, to:
$$
c_n = 1, \tr A = - c_{n - 1}, \determinant A = (-1)^n c_0
$$

\subsubsection{Twierdzenie Cayleya-Hamiltona}

Każda macierz kwadratowa nad $\mathbb{R}$ lub $\mathbb{C}$ jest
pierwiastkiem swojego wielomianu charakterystycznego.

$$
p_A(A) = 0
$$

Dla $p_A(t) = t^2 - 5A - 2I$, możemy bardzo łatwo określić wysokie potęgi $A$
lub jej odwrotność:
$$
p_A(A) = A^2 - 5A - 2I = 0 \Rightarrow A^2 = 5A + I \Rightarrow A^4 = (5A + I)^2
$$
$$
p_A(A) = 0 \Rightarrow A^2A^{-1} -5 - 2A^{-1} = 0 \Rightarrow A^{-1}
= \frac{A - 5I}{2}
$$

\subsubsection{Wielomian Minimalny}

Wielomian minimalny to wielomian $p_A(t)$, który spełnia warunek
$p_A(A) = 0$, oraz którego współczynnik przy najniższej potędze wynosi 1.
Każda macierz ma wielomian minimalny stopnia równego wielomianowi
charakterystycznemu. Macierz $A$ jest diagonalizowalna tylko wtedy, gdy $q_A$
ma pierwiastki jednorodne. Dzielnikami elementarnymi $f_i(\lambda)^{m_i}$
macierzy $A$ nazywamy potęgi nierozkładalnych wielomianów pojawiające się w
rozkładzie wielomianu minimalnego na czynniki.
$$
q_A(\lambda) = f_1(\lambda)^{m_1} f_2(\lambda)^{m_2}\dots f_k(\lambda)^{m_k}
$$

\subsection{Macierze podobne}

Macierze kwadratowe $A$ i $B$ nazywamy podobnymi $A \sim B$, jeśli istnieje $P$,
takie, że:
$$
PAP^{-1} = B
$$
Intuicyjnie, macierze są podobne jeśli da się je sprowadzić do siebie przez
jedną macierz. Jeśli $A \sim B$ to $p_A(t) = p_B(t)$.

\subsection{Diagonalizacja}

Macierz można zdiagonalizować tylko jeśli:
\begin{itemize}
  \item Jej wielomian charakterystyczny rozkłada się na iloczyn
    czynników liniowych
  \item Krotność geometryczna każdej wartości własnej odpowiada jej
    krotności algebraicznej
\end{itemize}

Macierze można jednocześnie diagonalizować, tylko jeśli ze sobą komutują:
$AB = BA$. Jednoczesna diagonalizacja oznacza istnienie jednej bazy
ortonormalnej, w której obie macierze przyjmują postać diagonalną.
Wektory tej bazy są więc wektorami własnymi zarówno A, jak i B.

\section{Klatka Jordana}

Macierz rozmiaru $k \times k$, postaci:
$$
J_k(\lambda) =
\begin{bmatrix}
  \lambda & 1 & 0 & \hdots & 0 \\
  0 & \lambda & 1 & \hdots & 0 \\
  \vdots & & \ddots & \ddots & \vdots \\
  0 & 0 & 0 & 0 & \lambda
\end{bmatrix}
$$
nazywamy klatką Jordana. Każda macierz kwadratowa $A$ o elementach z
$\mathbb{C}$ jest podobna do macierzy klatkowo-diagonalnej. Macierz klatkowo
diagonalną można zdefiniować jako macierz diagonalną, z kolejnymi
$J_1, \dots J_k$ na diagonali.

\end{document}
